% Options for packages loaded elsewhere
\PassOptionsToPackage{unicode}{hyperref}
\PassOptionsToPackage{hyphens}{url}
%
\documentclass[
]{article}
\usepackage{amsmath,amssymb}
\usepackage{iftex}
\ifPDFTeX
  \usepackage[T1]{fontenc}
  \usepackage[utf8]{inputenc}
  \usepackage{textcomp} % provide euro and other symbols
\else % if luatex or xetex
  \usepackage{unicode-math} % this also loads fontspec
  \defaultfontfeatures{Scale=MatchLowercase}
  \defaultfontfeatures[\rmfamily]{Ligatures=TeX,Scale=1}
\fi
\usepackage{lmodern}
\ifPDFTeX\else
  % xetex/luatex font selection
\fi
% Use upquote if available, for straight quotes in verbatim environments
\IfFileExists{upquote.sty}{\usepackage{upquote}}{}
\IfFileExists{microtype.sty}{% use microtype if available
  \usepackage[]{microtype}
  \UseMicrotypeSet[protrusion]{basicmath} % disable protrusion for tt fonts
}{}
\makeatletter
\@ifundefined{KOMAClassName}{% if non-KOMA class
  \IfFileExists{parskip.sty}{%
    \usepackage{parskip}
  }{% else
    \setlength{\parindent}{0pt}
    \setlength{\parskip}{6pt plus 2pt minus 1pt}}
}{% if KOMA class
  \KOMAoptions{parskip=half}}
\makeatother
\usepackage{xcolor}
\usepackage[margin=1in]{geometry}
\usepackage{color}
\usepackage{fancyvrb}
\newcommand{\VerbBar}{|}
\newcommand{\VERB}{\Verb[commandchars=\\\{\}]}
\DefineVerbatimEnvironment{Highlighting}{Verbatim}{commandchars=\\\{\}}
% Add ',fontsize=\small' for more characters per line
\usepackage{framed}
\definecolor{shadecolor}{RGB}{248,248,248}
\newenvironment{Shaded}{\begin{snugshade}}{\end{snugshade}}
\newcommand{\AlertTok}[1]{\textcolor[rgb]{0.94,0.16,0.16}{#1}}
\newcommand{\AnnotationTok}[1]{\textcolor[rgb]{0.56,0.35,0.01}{\textbf{\textit{#1}}}}
\newcommand{\AttributeTok}[1]{\textcolor[rgb]{0.13,0.29,0.53}{#1}}
\newcommand{\BaseNTok}[1]{\textcolor[rgb]{0.00,0.00,0.81}{#1}}
\newcommand{\BuiltInTok}[1]{#1}
\newcommand{\CharTok}[1]{\textcolor[rgb]{0.31,0.60,0.02}{#1}}
\newcommand{\CommentTok}[1]{\textcolor[rgb]{0.56,0.35,0.01}{\textit{#1}}}
\newcommand{\CommentVarTok}[1]{\textcolor[rgb]{0.56,0.35,0.01}{\textbf{\textit{#1}}}}
\newcommand{\ConstantTok}[1]{\textcolor[rgb]{0.56,0.35,0.01}{#1}}
\newcommand{\ControlFlowTok}[1]{\textcolor[rgb]{0.13,0.29,0.53}{\textbf{#1}}}
\newcommand{\DataTypeTok}[1]{\textcolor[rgb]{0.13,0.29,0.53}{#1}}
\newcommand{\DecValTok}[1]{\textcolor[rgb]{0.00,0.00,0.81}{#1}}
\newcommand{\DocumentationTok}[1]{\textcolor[rgb]{0.56,0.35,0.01}{\textbf{\textit{#1}}}}
\newcommand{\ErrorTok}[1]{\textcolor[rgb]{0.64,0.00,0.00}{\textbf{#1}}}
\newcommand{\ExtensionTok}[1]{#1}
\newcommand{\FloatTok}[1]{\textcolor[rgb]{0.00,0.00,0.81}{#1}}
\newcommand{\FunctionTok}[1]{\textcolor[rgb]{0.13,0.29,0.53}{\textbf{#1}}}
\newcommand{\ImportTok}[1]{#1}
\newcommand{\InformationTok}[1]{\textcolor[rgb]{0.56,0.35,0.01}{\textbf{\textit{#1}}}}
\newcommand{\KeywordTok}[1]{\textcolor[rgb]{0.13,0.29,0.53}{\textbf{#1}}}
\newcommand{\NormalTok}[1]{#1}
\newcommand{\OperatorTok}[1]{\textcolor[rgb]{0.81,0.36,0.00}{\textbf{#1}}}
\newcommand{\OtherTok}[1]{\textcolor[rgb]{0.56,0.35,0.01}{#1}}
\newcommand{\PreprocessorTok}[1]{\textcolor[rgb]{0.56,0.35,0.01}{\textit{#1}}}
\newcommand{\RegionMarkerTok}[1]{#1}
\newcommand{\SpecialCharTok}[1]{\textcolor[rgb]{0.81,0.36,0.00}{\textbf{#1}}}
\newcommand{\SpecialStringTok}[1]{\textcolor[rgb]{0.31,0.60,0.02}{#1}}
\newcommand{\StringTok}[1]{\textcolor[rgb]{0.31,0.60,0.02}{#1}}
\newcommand{\VariableTok}[1]{\textcolor[rgb]{0.00,0.00,0.00}{#1}}
\newcommand{\VerbatimStringTok}[1]{\textcolor[rgb]{0.31,0.60,0.02}{#1}}
\newcommand{\WarningTok}[1]{\textcolor[rgb]{0.56,0.35,0.01}{\textbf{\textit{#1}}}}
\usepackage{longtable,booktabs,array}
\usepackage{calc} % for calculating minipage widths
% Correct order of tables after \paragraph or \subparagraph
\usepackage{etoolbox}
\makeatletter
\patchcmd\longtable{\par}{\if@noskipsec\mbox{}\fi\par}{}{}
\makeatother
% Allow footnotes in longtable head/foot
\IfFileExists{footnotehyper.sty}{\usepackage{footnotehyper}}{\usepackage{footnote}}
\makesavenoteenv{longtable}
\usepackage{graphicx}
\makeatletter
\newsavebox\pandoc@box
\newcommand*\pandocbounded[1]{% scales image to fit in text height/width
  \sbox\pandoc@box{#1}%
  \Gscale@div\@tempa{\textheight}{\dimexpr\ht\pandoc@box+\dp\pandoc@box\relax}%
  \Gscale@div\@tempb{\linewidth}{\wd\pandoc@box}%
  \ifdim\@tempb\p@<\@tempa\p@\let\@tempa\@tempb\fi% select the smaller of both
  \ifdim\@tempa\p@<\p@\scalebox{\@tempa}{\usebox\pandoc@box}%
  \else\usebox{\pandoc@box}%
  \fi%
}
% Set default figure placement to htbp
\def\fps@figure{htbp}
\makeatother
\setlength{\emergencystretch}{3em} % prevent overfull lines
\providecommand{\tightlist}{%
  \setlength{\itemsep}{0pt}\setlength{\parskip}{0pt}}
\setcounter{secnumdepth}{-\maxdimen} % remove section numbering
\usepackage{xcolor}
\makeatletter
\makeatother
\usepackage{float}
\renewcommand{\CommentTok}[1]{\textcolor[rgb]{0.60,0.20,0.20}{#1}}
\usepackage{bookmark}
\IfFileExists{xurl.sty}{\usepackage{xurl}}{} % add URL line breaks if available
\urlstyle{same}
\hypersetup{
  pdftitle={SC306 Homework 5},
  pdfauthor={Othar Zaldastani II},
  hidelinks,
  pdfcreator={LaTeX via pandoc}}

\title{SC306 Homework 5}
\author{Othar Zaldastani II}
\date{March 17, 2026}

\begin{document}
\maketitle

\section{Part 1: AR and PAR}\label{part-1-ar-and-par}

\paragraph{A)}\label{a}

\[AR \ \% = \frac{R_e - R_u}{R_e} \times 100\]

\begin{verbatim}
## [1] 77.1
\end{verbatim}

The attributable risk percent of low birth weight from Cocaine use is
77.1\%.

\paragraph{B)}\label{b}

This means approximately 77\% of the low birth weight risk for infants
born to mothers who used Cocaine while pregnant is attributable to their
use of Cocaine.

\paragraph{C)}\label{c}

\[PAR \ \% = \frac{P_e(R_e - R_u)}{P_e R_e + (1-P_e)R_u} \times 100\]

\begin{verbatim}
## [1] 6.3
\end{verbatim}

The population attributable risk percent of low birth weight from
Cocaine use is 6.3\%.

\paragraph{D)}\label{d}

Approximately 6.3\% of the low birth weight risk for all infants born in
California is attributable to the use of Cocaine while pregnant.

\pagebreak

\section{Part 2: Study Designs}\label{part-2-study-designs}

\subsection{Abstract 1}\label{abstract-1}

The researchers used a randomized control trial (RCT) with a six-month
follow up.

Effects of reduction in screen use/exposure via the use of a modified
classroom curriculum.

Childhood obesity related measurements such as BMI and bodily
dimensions/ratios.

The primary measures used to quantify disease were obesity-related
measures such as BMI and body measurements. Since a condition like
obesity follows more of a sliding scale of severity, changes in bodily
characteristics were used as opposed to number of participants
over/under a threshold.

The measures used to quantify exposure vs.~outcome relationship was
difference in means (difference in differences) with 95\% confidence
intervals.

\subsection{Abstract 2}\label{abstract-2}

The researchers used a cohort study with continuous, indefinite follow
ups to conduct this experiment. The cohort was followed from the birth
period of 1915-1929 to 1995 and mortality causes were assessed.

The exposure investigated was fetal growth rate (controlling for size at
birth).

The disease investigated was mortality from Ischaemic Heart Disease (and
other causes) during the lives of the participants.

The measures used to quantify disease were mortality rates from
Ischaemic Heart Disease and other Cardiovascular diseases.

The measure used to quantify the relationship between exposure and
outcome was incidence rate ratios relative to birth weight with 95\%
confidence intervals. An additional set of results was published
controlling for socioeconomic factors at birth and adulthood.

\subsection{Abstract 3}\label{abstract-3}

The researchers used a case-control study. Women with premenopausal
breast cancer were compared to demographically-matched controls. Prior
history of electric blanket use was assessed retrospectively.

The exposure investigated was electric blanket use (as a proxy for
increased background exposure to 60 hz electromagnetic fields). Data on
blanket use included time, duration, and frequency of use over the
previous decade.

The disease investigated was premenopausal breast cancer.

Since this is a case-control study, disease frequency is not a
meaningful statistic. Numbers of cases and controls were chosen at the
start of the study.

The measure used to quantify relationship between electric blanket use
and premenopausal breast cancer are odds ratios with 95\% confidence
intervals. The odds ratios were adjusted for age, education, and other
possible risk factors/confounding variables.

\pagebreak

\section{Part 3: A Case-Control
Study}\label{part-3-a-case-control-study}

\paragraph{A)}\label{a-1}

Since cervical cancer is relatively uncommon, it makes sense to do a
case-control study compared to other options such as an RCT or cohort
study. An RCT would be challenging and extremely unethical, while a
cohort study would require an immense number of participants and time in
order to have a sufficient number of female participants develop
cervical cancer. A case-control study is best because the disease takes
a long time to incubate, and the retrospective information needed to
categorize participants is not very granular (number of sexual partners,
sex worker visitation, etc.) so there is a lower likelihood of receiving
false information.

\paragraph{B)}\label{b-1}

Some potential problems include recall bias, reporting bias, selection
bias, and confounding issues. Cases and their husbands may put more
effort into recalling past information compared to controls. Contact
with sex workers is stigmatized and some husbands may lie about their
history. The selection of individuals from hospitals also poses a
potential problem because hospitalized individuals are not a
representative segment of the population. Also, some potential
confounding variables include exposure to carcinogens such as cigarette
use, socioeconomic status, occupation or occupational hazard and
exposure, and the wives' number of sexual partners prior to (or during)
marriage.

\paragraph{C)}\label{c-1}

2x2 table of cervical cancer data

\begin{longtable}[]{@{}lrr@{}}
\toprule\noalign{}
& cancer & No cancer \\
\midrule\noalign{}
\endhead
\bottomrule\noalign{}
\endlastfoot
CSW & 179 & 518 \\
No CSW & 42 & 267 \\
\end{longtable}

\paragraph{D)}\label{d-1}

Women with husbands who had previously visited sex workers had higher
odds of cervical cancer compared to women with husbands who had reported
never having visited a sex worker.

\paragraph{E)}\label{e}

The prevalence of cervical cancer among study participants is 22\%.
\[\frac{\text{\# of Cases in Study}} {\text{\# of Participants in Study}} = \frac{179+42}{179 + 518 + 42 + 267} = 0.220\]

\begin{verbatim}
## [1] 0.22
\end{verbatim}

\paragraph{F)}\label{f}

No, because in a case-cohort study, participants are not randomly
selected but, so the study population has a significantly higher rate of
the condition (here, cervical cancer) than the general population.
Investigators intentionally set the number of cases and controls in the
study.

\paragraph{G)}\label{g}

The 95\% confidence interval for the odds ratio is \([1.53 : 3.2]\)
centered at 2.19.

\begin{verbatim}
## estimate    lower    upper 
##     2.19     1.53     3.20
\end{verbatim}

\paragraph{H)}\label{h}

The odds ratio of 2.19 with a 95\% confidence interval of
\([1.53 : 3.2]\) means that the odds of developing invasive cervical
cancer among women with husbands who had previously visited sex workers
is estimated to be between 1.5 and 3.2 times the odds of the general
population with husbands who did not visit sex workers.

\pagebreak

\section{Code Appendix}\label{code-appendix}

\begin{Shaded}
\begin{Highlighting}[numbers=left,,]
\NormalTok{knitr}\SpecialCharTok{::}\NormalTok{opts\_chunk}\SpecialCharTok{$}\FunctionTok{set}\NormalTok{(          }\CommentTok{\# setup chunk}
  \AttributeTok{echo =} \ConstantTok{FALSE}\NormalTok{,                 }
  \AttributeTok{warning =} \ConstantTok{FALSE}\NormalTok{,              }
  \AttributeTok{message =} \ConstantTok{FALSE}\NormalTok{,              }
  \AttributeTok{cache =} \ConstantTok{FALSE}\NormalTok{,                }
  \AttributeTok{fig.align =} \StringTok{"center"}\NormalTok{,         }
  \AttributeTok{attr.source =} \StringTok{\textquotesingle{}.numberLines\textquotesingle{}}\NormalTok{,}
  \AttributeTok{cache =}\NormalTok{ T)}

\FunctionTok{library}\NormalTok{(tidyverse)              }\CommentTok{\# loading packages}
\FunctionTok{library}\NormalTok{(dplyr)}
\FunctionTok{library}\NormalTok{(knitr)}
\FunctionTok{library}\NormalTok{(epitools)}

\FunctionTok{set.seed}\NormalTok{(}\DecValTok{24}\NormalTok{)                    }\CommentTok{\# seed for reproducibility}

\NormalTok{kable2 }\OtherTok{\textless{}{-}} \ControlFlowTok{function}\NormalTok{(x, ...) \{    }\CommentTok{\# for debugging table creation}
\NormalTok{  kableExtra}\SpecialCharTok{::}\FunctionTok{kable\_styling}\NormalTok{(}
\NormalTok{    knitr}\SpecialCharTok{::}\FunctionTok{kable}\NormalTok{(x, ...))\}}
\CommentTok{\# Part 1}
\CommentTok{\# Question A}
\NormalTok{Re }\OtherTok{=}\NormalTok{ (}\DecValTok{110}\NormalTok{)}\SpecialCharTok{/}\NormalTok{(}\DecValTok{110}\SpecialCharTok{+}\DecValTok{289}\NormalTok{)            }\CommentTok{\# Risk among exposed}
\NormalTok{Ru }\OtherTok{=}\NormalTok{ (}\DecValTok{1073}\NormalTok{)}\SpecialCharTok{/}\NormalTok{(}\DecValTok{1073}\SpecialCharTok{+}\DecValTok{15896}\NormalTok{)        }\CommentTok{\# Risk among unexposed}
\NormalTok{AR }\OtherTok{=}\NormalTok{ (Re }\SpecialCharTok{{-}}\NormalTok{ Ru)}\SpecialCharTok{/}\NormalTok{(Re) }\SpecialCharTok{*} \DecValTok{100}       \CommentTok{\# Attributable risk}
\FunctionTok{round}\NormalTok{(AR, }\DecValTok{1}\NormalTok{)}

\CommentTok{\# Question C}
\NormalTok{Pe }\OtherTok{=} \FloatTok{0.02}                       \CommentTok{\# Exposure prevalence}
\NormalTok{PAR }\OtherTok{=}\NormalTok{ (Pe }\SpecialCharTok{*}\NormalTok{ (Re }\SpecialCharTok{{-}}\NormalTok{ Ru)) }\SpecialCharTok{/}\NormalTok{ (Pe }\SpecialCharTok{*}\NormalTok{ Re }\SpecialCharTok{+}\NormalTok{ (}\DecValTok{1}\SpecialCharTok{{-}}\NormalTok{Pe) }\SpecialCharTok{*}\NormalTok{ Ru) }\SpecialCharTok{*} \DecValTok{100}
\FunctionTok{round}\NormalTok{(PAR, }\DecValTok{2}\NormalTok{)}

\CommentTok{\# Part 3}
\CommentTok{\# Question C}
\NormalTok{table.cancer }\OtherTok{=}                  \CommentTok{\# Creation of 2x2 table}
  \FunctionTok{matrix}\NormalTok{(}\FunctionTok{c}\NormalTok{(}\DecValTok{179}\NormalTok{, }\DecValTok{518}\NormalTok{, }\DecValTok{42}\NormalTok{, }\DecValTok{267}\NormalTok{),}\AttributeTok{nrow=}\DecValTok{2}\NormalTok{, }\AttributeTok{byrow =} \ConstantTok{TRUE}\NormalTok{)}
\FunctionTok{colnames}\NormalTok{(table.cancer) }\OtherTok{=} \FunctionTok{c}\NormalTok{(}\StringTok{\textquotesingle{}cancer\textquotesingle{}}\NormalTok{, }\StringTok{\textquotesingle{}No cancer\textquotesingle{}}\NormalTok{)}
\FunctionTok{rownames}\NormalTok{(table.cancer) }\OtherTok{=} \FunctionTok{c}\NormalTok{(}\StringTok{\textquotesingle{}CSW\textquotesingle{}}\NormalTok{, }\StringTok{\textquotesingle{}No CSW\textquotesingle{}}\NormalTok{)}
\FunctionTok{kable}\NormalTok{(table.cancer)             }\CommentTok{\# Printing 2x2 table}

\CommentTok{\# Question E}
\FunctionTok{round}\NormalTok{(((}\DecValTok{179}\SpecialCharTok{+}\DecValTok{42}\NormalTok{)}\SpecialCharTok{/}\NormalTok{(}\DecValTok{179} \SpecialCharTok{+} \DecValTok{518} \SpecialCharTok{+} \DecValTok{42} \SpecialCharTok{+} \DecValTok{267}\NormalTok{)), }\DecValTok{3}\NormalTok{)}
\CommentTok{\# Question G}
\NormalTok{or }\OtherTok{=} \FunctionTok{oddsratio}\NormalTok{(table.cancer, }\AttributeTok{rev =} \StringTok{"b"}\NormalTok{)}
\FunctionTok{round}\NormalTok{(or}\SpecialCharTok{$}\NormalTok{measure[}\StringTok{"CSW"}\NormalTok{, ], }\DecValTok{2}\NormalTok{)}
\end{Highlighting}
\end{Shaded}

\vfill

Attributions: I have used some template/styling code from GitHub user
alexpghayes.

\end{document}
